\chapter{Syphilis Testing}

\section*{What Is This Test?}
Syphilis testing detects \textit{Treponema pallidum} subspecies \textit{pallidum}, the causative agent of syphilis, an obligate human pathogen that cannot be cultured in vitro. Syphilis progresses through distinct clinical stages: primary (painless chancre at the inoculation site 10--90 days post-exposure), secondary (disseminated rash including palmar/plantar involvement, mucocutaneous lesions, condyloma lata, lymphadenopathy 4--10 weeks after the chancre), early latent (asymptomatic, <1 year from acquisition), late latent (asymptomatic, >1 year), and tertiary (cardiovascular syphilis, neurosyphilis, or gummatous disease occurring years to decades after initial infection). Congenital syphilis results from transplacental transmission, causing stillbirth, neonatal death, or lifelong sequelae. Testing employs two algorithm types: the traditional algorithm (nontreponemal screening test followed by treponemal confirmatory test) and the reverse algorithm (automated treponemal immunoassay as the screening test, followed by nontreponemal test, with discordant results resolved by a second treponemal test). Nontreponemal tests (RPR, VDRL) measure IgG and IgM antibodies against cardiolipin-lecithin-cholesterol antigen released from damaged host cells; they are semi-quantitative and titers correlate with disease activity. Treponemal tests (TPPA/TPHA, FTA-ABS, EIA/CIA) detect antibodies specific to \textit{T. pallidum} antigens; they remain positive for life in most individuals.

\section*{Alternative Names}
\begin{itemize}
  \item Syphilis Serology
  \item RPR
  \item VDRL
  \item Rapid Plasma Reagin
  \item Treponemal Test
  \item TPPA
  \item TPHA
  \item FTA-ABS
  \item Syphilis IgG/IgM
  \item Syphilis EIA
  \item Venereal Disease Research Laboratory Test
  \item Syphilis Screening Cascade
\end{itemize}
\section*{Principle}
Nontreponemal tests (RPR, VDRL): The RPR (rapid plasma reagin) test uses charcoal-cardiolipin particles for macroscopic agglutination reading. The VDRL (Venereal Disease Research Laboratory) test is a microscopic flocculation test, used primarily for CSF (VDRL-CSF for neurosyphilis diagnosis). Both detect anti-cardiolipin antibodies. Results are reported semi-quantitatively as titers (1:1, 1:2, 1:4, 1:8, 1:16, 1:32, 1:64, etc.). Sensitivity varies by stage: 78--86\% for primary, 100\% for secondary, 95--100\% for early latent, 71--73\% for late latent/tertiary. Specificity 98--99\%, but false positives occur in pregnancy, autoimmune disease, HIV, acute febrile illness, and injection drug use. A fourfold change in titer (e.g., from 1:16 to 1:4, or 1:4 to 1:16) is considered clinically significant. Treponemal tests (EIA/CIA, TPPA, FTA-ABS): Enzyme/chemiluminescent immunoassays detect antibodies to recombinant \textit{T. pallidum} antigens (TpN15, TpN17, TpN47). The \textit{T. pallidum} particle agglutination assay (TPPA) uses gelatin particles sensitized with \textit{T. pallidum} antigens. The fluorescent treponemal antibody absorption test (FTA-ABS) uses indirect immunofluorescence. Treponemal tests are highly sensitive (>85\% in primary syphilis, 100\% in secondary and latent) and specific (96--99\%) but remain positive for life regardless of treatment and cannot distinguish active from treated infection. PCR: Detects \textit{T. pallidum} DNA (targeting \textit{tp47} or \textit{polA} genes) from primary chancre exudate, CSF, or tissue. Sensitivity 80--95\%, specificity >95\%. Particularly useful for primary syphilis (before seroconversion), congenital syphilis, and neurosyphilis. Darkfield microscopy: Direct visualization of motile spirochetes from moist primary or secondary lesions; sensitivity 75--100\% depending on operator expertise; now rarely available outside specialized STD clinics.

\section*{History}
Syphilis was first documented in Europe in 1495. \textit{Treponema pallidum} was identified by Fritz Schaudinn and Erich Hoffmann in 1905. The Wassermann complement fixation test (the first nontreponemal test) was developed in 1906. The VDRL test was developed in 1946, and the RPR in the 1960s. The FTA-ABS treponemal test was introduced in the 1960s. Penicillin was shown to be curative for syphilis in the 1940s; remarkably, \textit{T. pallidum} has remained universally susceptible to penicillin. The HIV epidemic in the 1980s complicated syphilis diagnosis due to atypical serologic responses in HIV co-infected individuals. The reverse sequence screening algorithm (automated treponemal EIA/CLIA first) gained popularity in the 2000s as high-throughput immunoassay platforms became standard in clinical laboratories. The resurgence of syphilis in the 2000s--2010s, including increasing rates of congenital syphilis (reaching a 20-year high in 2019 and continuing to increase), has prompted CDC and USPSTF to reinforce screening recommendations for pregnant women and high-risk populations.

\section*{Why This Test Is Ordered}
\begin{itemize}
  \item Screening for syphilis in individuals at increased risk: MSM (annually or every 3--6 months depending on risk behaviors), individuals with HIV, commercial sex workers, individuals with multiple sexual partners, and individuals diagnosed with another STI
  \item Universal prenatal screening: All pregnant women should be screened for syphilis at the first prenatal visit; repeat screening at 28--32 weeks of gestation and at delivery is recommended for women at high risk or in communities with high syphilis prevalence
  \item Diagnosis of genital, anal, or oral ulcerative lesions to distinguish syphilitic chancre (typically painless, indurated, clean-based) from HSV (painful, multiple, vesicular), chancroid (\textit{Haemophilus ducreyi}, painful, ragged, purulent), LGV, and granuloma inguinale
  \item Evaluation of patients with unexplained generalized rash (including palms and soles), mucocutaneous lesions, generalized lymphadenopathy, or alopecia---the classic manifestations of secondary syphilis
  \item Evaluation for neurosyphilis in patients with neurologic symptoms (stroke, altered mental status, visual changes, hearing loss, tabes dorsalis, general paresis) and reactive syphilis serology
  \item Evaluation of patients with unexplained ocular inflammation (uveitis, retinitis, optic neuritis), which may be the presenting manifestation of ocular syphilis at any stage
  \item Screening of organ and blood donors for syphilis, as required by FDA and transplant regulations
  \item Monitoring of treatment response: serial RPR titers at 3, 6, 12, and 24 months after treatment. A fourfold decline in RPR titer (e.g., 1:32 to 1:8) indicates adequate treatment response
\end{itemize}

\section*{Primary vs Secondary Test}
Syphilis serology is the primary screening test. Nontreponemal tests (RPR/VDRL) with treponemal confirmation, or the reverse algorithm (treponemal EIA first), are primary screening strategies. Darkfield microscopy and PCR are primary diagnostic tests when active lesions are present, with the advantage of diagnosing primary syphilis before seroconversion occurs.

\section*{How This Test Is Performed}
Venous blood collected in an SST tube for serology; 5 mL is sufficient. For CSF VDRL: CSF collected by lumbar puncture in a sterile tube; at least 0.5 mL is required for VDRL. For PCR: Lesion exudate collected by gently abrading the lesion base with a sterile Dacron swab after cleaning with saline (NOT with alcohol or antiseptic); place swab in viral transport medium or manufacturer-specific transport medium. For darkfield microscopy: Lesion must be cleaned with saline, abraded to produce serous exudate (blood should be avoided as it obscures the field), and examined immediately (within 20 minutes) by an experienced microscopist. Serum for serology can be refrigerated at 2--8\textdegree{}C for up to 7 days or frozen for extended storage.

\section*{What Preparation Is Needed}
No special preparation for venipuncture. Patients should avoid applying topical medications, creams, or antibiotics to lesions before specimen collection for PCR or darkfield microscopy. The patient should not urinate immediately before genital lesion sampling. Alcohol or antiseptic cleaning of the lesion before darkfield or PCR collection should be avoided as it may kill spirochetes or degrade nucleic acid; saline irrigation only.

\section*{Contraindications and Precautions}
\begin{itemize}
  \item Nontreponemal tests (RPR, VDRL) are subject to the prozone phenomenon: in secondary syphilis, when antibody titers are extremely high (>1:512), the undiluted serum may produce a false-negative result because antigen-antibody complexes form in antibody excess and fail to agglutinate. The laboratory must dilute the specimen to resolve the prozone. False-positive nontreponemal tests (biologic false positive, BFP) occur in 1--2\% of the general population and more commonly in: pregnancy, autoimmune disease (SLE, antiphospholipid syndrome), HIV, injection drug use, acute febrile illnesses, and older adults (>70 years). BFPs are typically low titer (<1:8) and are not confirmed by treponemal tests. Treponemal tests remain positive for life in most treated patients (the "serologic scar") and cannot distinguish active from treated infection. Discordant results in the reverse algorithm (EIA+/RPR-) require resolution with a second treponemal test (TPPA). The treponemal EIA/CLIA may remain positive after treatment and must NOT be used to monitor treatment response. The Jarisch-Herxheimer reaction (fever, myalgia, headache, exacerbation of rash) occurs in up to 30--50\% of patients within 24 hours of treatment, most commonly in early syphilis. It is self-limited and does not contraindicate continued therapy.
\end{itemize}
\section*{Risks}
Venipuncture and CSF collection carry standard procedural risks. The major clinical risk is failure to diagnose syphilis in its early (most infectious) stages. Undiagnosed and untreated syphilis in pregnancy leads to congenital syphilis in >70\% of cases with high rates of stillbirth, neonatal death, and lifelong morbidity. All reactive serology must be promptly communicated and acted upon.

\section*{What Happens After the Test}
Serology results available in 1--24 hours. PCR results in 24--72 hours. Darkfield microscopy results immediately. A reactive RPR or treponemal EIA requires further testing per the algorithm. Confirmed syphilis requires prompt treatment with benzathine penicillin G (2.4 million units IM as a single dose for early syphilis; 2.4 million units IM weekly x 3 doses for late latent syphilis or unknown duration). Patients must be counseled to abstain from sexual activity until lesions are completely healed (7 days after treatment). Sexual partners from the relevant time period (3 months + duration of symptoms for primary; 6 months + duration for secondary; 1 year for early latent) must be evaluated and treated presumptively. All cases must be reported to public health.

\section*{Understanding Results}

\subsection*{Below Normal}
\begin{itemize}
  \item \textbf{RPR nonreactive:} No detectable anti-cardiolipin antibodies. May indicate: no syphilis infection; very early primary syphilis (chancre present <1--2 weeks, before seroconversion---RPR sensitivity in primary syphilis is 78--86\%); treated syphilis with seroreversion; late latent/tertiary syphilis with low titers (prozone resolved); or immunosuppression with impaired antibody response. In a patient with a genital ulcer and negative RPR, darkfield microscopy or PCR of the lesion should be performed
  \item \textbf{Treponemal EIA/CLIA negative:} No treponemal antibodies detected. Essentially excludes syphilis in immunocompetent individuals, except in very early primary infection (first 2 weeks after chancre appearance)
  \item \textbf{Twofold RPR titer decline after treatment (e.g., 1:16 to 1:8):} Consistent with treatment response; monitoring should continue for a fourfold decline
\end{itemize}

\subsection*{Above Normal}
\begin{itemize}
  \item \textbf{RPR reactive (reported as titer, e.g., 1:32) with positive treponemal test (EIA or TPPA):} Confirms syphilis. A fourfold titer increase (e.g., 1:4 to 1:16) in a previously treated patient with known baseline titer indicates reinfection or treatment failure. In a new patient, the RPR titer combined with the clinical stage guides treatment (early vs. late)
  \item \textbf{RPR reactive with negative treponemal test (Traditional algorithm with RPR first; in reverse algorithm, EIA+/RPR+/TPPA+ confirms syphilis):} This constitutes a biologic false positive (BFP) and does NOT represent syphilis. The patient should be evaluated for causes of BFP (pregnancy, autoimmune disease, HIV, acute infection, advanced age). Follow-up testing in 4--6 weeks may be considered
  \item \textbf{Treponemal EIA/CLIA reactive, RPR nonreactive, TPPA reactive (reverse algorithm):} Indicates past (treated or spontaneously resolved) syphilis, very early primary syphilis (RPR not yet reactive), late latent syphilis with low titer, or a false-positive EIA with true-positive TPPA. A detailed clinical and treatment history is essential. If prior treatment is documented and adequate, no further treatment is required. If no prior treatment is documented or it was inadequate, treat for late latent syphilis
  \item \textbf{CSF VDRL reactive:} Diagnostic of neurosyphilis (highly specific). CSF WBC >5 cells/$\mu$L and/or protein >45 mg/dL supports the diagnosis. A nonreactive CSF VDRL does NOT exclude neurosyphilis (sensitivity 30--70\%). CSF treponemal test (FTA-ABS) is highly sensitive (>95\%) and a nonreactive CSF FTA-ABS makes neurosyphilis very unlikely
  \item \textbf{PCR positive for \textit{T. pallidum} DNA:} Confirms active syphilis. Useful in primary syphilis (chancre), congenital syphilis (placenta, amniotic fluid, neonatal CSF), and secondary syphilis (lesion exudate). Not approved by FDA but available through reference laboratories
\end{itemize}

\section*{Normal Results}
\textbf{RPR nonreactive, Treponemal test (EIA, TPPA) nonreactive} (both algorithms). No evidence of treponemal infection. CSF VDRL: \textbf{Nonreactive}. CSF WBC: \textbf{0--5 cells/$\mu$L}. CSF protein: \textbf{15--45 mg/dL}. Darkfield: \textbf{No spirochetes visualized}.

\section*{Follow-up Tests}
\begin{itemize}
  \item \textbf{RPR titer monitoring:} At 3, 6, 12, and 24 months after treatment. A fourfold or greater decline (e.g., 1:32 to 1:8) indicates adequate treatment response. Failure of the RPR titer to decline fourfold by 6--12 months (serofast state), or a fourfold or greater increase in titer from a prior result, requires evaluation for treatment failure or reinfection
  \item \textbf{CSF examination (cell count, protein, VDRL):} Indicated for: neurologic or ocular symptoms/signs; evidence of tertiary syphilis; treatment failure; HIV infection with CD4 <350 and/or RPR $\geq$1:32 (some guidelines recommend CSF examination in this scenario)
  \item \textbf{Treponemal confirmatory test (TPPA or FTA-ABS):} For RPR-reactive patients to distinguish true syphilis from biologic false positive
  \item \textbf{HIV testing:} All patients diagnosed with syphilis should be tested for HIV, as syphilis increases HIV acquisition and transmission risk 3- to 5-fold
  \item \textbf{Other STI screening:} Gonorrhea, chlamydia (NAAT at exposed sites), and hepatitis B/C testing per CDC recommendations
  \item \textbf{Darkfield microscopy or PCR of suspicious lesions:} In patients with suspected primary or secondary syphilis and negative serology to establish the diagnosis before seroconversion
\end{itemize}

\section*{References}
\begin{enumerate}
  \item Workowski KA, Bachmann LH, Chan PA, et al. Sexually transmitted infections treatment guidelines, 2021. MMWR Recomm Rep. 2021;70(4):1-187.
  \item CDC. Syphilis --- CDC fact sheet (detailed). Centers for Disease Control and Prevention; 2022.
  \item Janier M, Hegyi V, Dupin N, et al. 2014 European guideline on the management of syphilis. J Eur Acad Dermatol Venereol. 2014;28(12):1581-1593.
  \item Ghanem KG, Ram S, Rice PA. The modern epidemic of syphilis. N Engl J Med. 2020;382(9):845-854.
  \item Tuddenham S, Hamill MM, Ghanem KG. Diagnosis and treatment of sexually transmitted infections: a review. JAMA. 2022;327(2):161-172.
\end{enumerate}

\clearpage
\section*{Regulatory Status and Authorization Date}
This chapter describes a test category, not a specific commercial product. FDA, CDSCO, EU IVDR, and MHRA approval or registration dates therefore cannot be assigned without the assay manufacturer, product name, instrument, and intended use. For laboratory-developed tests, an individual agency approval date may not exist. See the Preface for the interpretation of regulatory status.

\clearpage
